\section{Worked metric tutorial}
\label{app:tutorial}

All five measures use the same survey distributions but summarize disagreement at
different scales. The worked example follows the calculation from one item to a country
profile and then to cross-country geometry; every value can be regenerated with
\texttt{scripts/make\_wave1\_assets.py}.

\subsection{Step 1: compare two response distributions}
For ordered categories $r=1,\ldots,R$, let $q_r$ be the human probability and $p_r$ the
model probability. Total variation ignores order,
\[
\mathrm{TVD}(p,q)=\frac12\sum_{r=1}^{R}|p_r-q_r|,
\]
whereas normalized one-dimensional Wasserstein distance measures how far probability
mass moves,
\[
W1(p,q)=\frac{1}{R-1}\sum_{r=1}^{R-1}|F_p(r)-F_q(r)|.
\]
Table~\ref{tab:tutorial-q106} works through China Q106 using the human distribution and
the mean of five model generations. For this item, GPT-5.6 has slightly lower TVD but
higher W1: it moves less total mass, but some of that mass travels farther along the
ordered scale.

\begin{table}[!htbp]
\centering
\caption{Worked item example: China, Q106. Model columns average the five generations.}
\label{tab:tutorial-q106}
\small
\begin{tabular}{rrrr}
\toprule
\rowcolor{tablehead} Score & Human $q_r$ & GPT-5.5 $p_r$ & GPT-5.6 Sol $p_r$ \\
\midrule
1 & 0.145 & 0.058 & 0.061 \\
2 & 0.072 & 0.040 & 0.031 \\
3 & 0.101 & 0.064 & 0.058 \\
4 & 0.029 & 0.074 & 0.055 \\
5 & 0.116 & 0.194 & 0.188 \\
6 & 0.145 & 0.114 & 0.104 \\
7 & 0.101 & 0.128 & 0.122 \\
8 & 0.087 & 0.140 & 0.160 \\
9 & 0.043 & 0.070 & 0.071 \\
10 & 0.159 & 0.118 & 0.150 \\
\midrule
W1 & -- & 0.0742 & 0.0913 \\
TVD & -- & 0.2292 & 0.2192 \\
\bottomrule
\end{tabular}
\end{table}
\begin{table}[!htbp]
\centering
\caption{Worked country-profile example for China. Directed scores are on $[0,1]$; $\sigma_k$ is the 64-country human SD.}
\label{tab:tutorial-crg}
\small
\begin{tabular}{lrrrrr}
\toprule
\rowcolor{tablehead} Domain & Human & GPT-5.5 & GPT-5.6 Sol & $\sigma_k$ & GPT-5.5 $z$ error \\
\midrule
Agency & 0.760 & 0.668 & 0.640 & 0.108 & -0.848 \\
Trust & 0.615 & 0.622 & 0.630 & 0.185 & 0.036 \\
Distribution & 0.494 & 0.433 & 0.405 & 0.162 & -0.380 \\
Responsibility & 0.426 & 0.473 & 0.440 & 0.105 & 0.448 \\
Immigration & 0.647 & 0.557 & 0.613 & 0.126 & -0.710 \\
Science & 0.856 & 0.787 & 0.754 & 0.080 & -0.862 \\
\midrule
CRG & -- & 0.621 & 0.736 & -- & -- \\
\bottomrule
\end{tabular}
\end{table}
\begin{table}[!htbp]
\centering
\caption{Four-country tutorial subset (China, Brazil, Nigeria, United States). This subset illustrates the arithmetic and is not a separate inferential result.}
\label{tab:tutorial-geometry}
\small
\begin{tabular}{llrrr}
\toprule
\rowcolor{tablehead} Country 1 & Country 2 & Human distance & GPT-5.5 & GPT-5.6 Sol \\
\midrule
Brazil & Nigeria & 0.448 & 0.130 & 0.155 \\
Brazil & United States & 0.323 & 0.329 & 0.311 \\
China & Brazil & 0.612 & 0.565 & 0.574 \\
China & Nigeria & 0.561 & 0.525 & 0.537 \\
China & United States & 0.366 & 0.324 & 0.353 \\
Nigeria & United States & 0.480 & 0.261 & 0.254 \\
\midrule
VDR & -- & 1.000 & 0.765 & 0.783 \\
CSR & -- & 1.000 & 0.486 & 0.543 \\
\bottomrule
\end{tabular}
\end{table}


\subsection{Step 2: turn each item into a directed country coordinate}
If higher raw values already point toward the named construct, the directed expected
score is
\[
s_{mcj}=\sum_r p_{mcjr}\frac{r-r_{\min}}{r_{\max}-r_{\min}}.
\]
For reverse-coded constructs, we use $1-s_{mcj}$. With one anchor per domain in the new
experiment, the six directed scores form
$A_{mc}=(A,T,D,M,I,S)$. The corresponding survey profile is $H_c$.

\subsection{Step 3: standardize the country-profile gap}
Let $\sigma_k$ be the across-country standard deviation of the human score for domain
$k$. Then
\[
\CRG_{mc}=\sqrt{\frac{1}{6}\sum_{k=1}^{6}
\left(\frac{A_{mck}-H_{ck}}{\sigma_k}\right)^2}.
\]
For example, each China row in Table~\ref{tab:tutorial-crg} first forms a standardized
error; squaring, averaging the six values, and taking the square root yields the displayed
CRG. CRG is zero only under exact agreement on all six coordinates.

\subsection{Step 4: compare cross-country geometry}
For every unordered pair $c<c'$, compute Euclidean distances
$d^H_{cc'}=\lVert H_c-H_{c'}\rVert_2$ and
$d^A_{m,cc'}=\lVert A_{mc}-A_{mc'}\rVert_2$. The dispersion ratio and structure
retention are
\[
\VDR_m=\frac{\overline d^A_m}{\overline d^H},
\qquad
\CSR_m=\rho_S(\mathbf d^A_m,\mathbf d^H).
\]
VDR asks whether the \emph{amount} of separation is preserved; CSR asks whether the
\emph{ordering} of which country pairs are near or far is preserved. The four-country
subset in Table~\ref{tab:tutorial-geometry} illustrates six pair distances. The reported
main estimates repeat exactly this calculation over all
$\binom{64}{2}=2{,}016$ pairs. The subset is pedagogical and is not a second analysis.

\begin{auditbox}{Reading the targets}
W1, TVD, and CRG target 0. VDR targets 1: values below 1 mean compression and values
above 1 mean expansion. CSR targets 1; 0 means no monotonic preservation of the human
pairwise ordering, and negative values indicate reversal.
\end{auditbox}
\FloatBarrier
